
\usepackage[a4paper, total={6in, 8in}]{geometry}
\usepackage{xcolor}
\usepackage{titlesec}
\usepackage{enumitem}
\setlist[itemize]{leftmargin=3em}
\usepackage{indentfirst}

\titleformat{\section}[block]{\color{black}\Large\bfseries\filcenter}{}{1em}{}
\usepackage{polski}
\usepackage{fontspec}
\usepackage{hyperref}
\usepackage[normalem]{ulem}
\usepackage{soul}
\usepackage{hyperref}
\usepackage{listings}
\usepackage{graphicx}

\usepackage{fancyhdr}

% Setup fancyhdr for all pages
\pagestyle{fancy}
\fancyhf{} % clear header and footer
\fancyfoot[C]{\date{\monthyearpl} - \pdfTitle - [\thepage]} % center footer with page number

% Make the plain style (used on title page) the same
\fancypagestyle{plain}{
	\fancyfoot[C]{\date{\monthyearpl} - \pdfTitle - [\thepage]} % center footer with page number
}


\hypersetup{
	colorlinks=true,
	linkcolor=blue,
	filecolor=magenta,      
	urlcolor=blue,
	pdftitle={\pdfTitle},
	pdfpagemode=FullScreen,
}
\usepackage{xcolor}

\colorlet{punct}{red!60!black}
\definecolor{background}{HTML}{EEEEEE}
\definecolor{delim}{RGB}{20,105,176}
\colorlet{numb}{magenta!60!black}
\lstdefinelanguage{json}{
	basicstyle=\normalfont\ttfamily,
	numbers=left,
	numberstyle=\scriptsize,
	stepnumber=1,
	numbersep=8pt,
	showstringspaces=false,
	breaklines=true,
	frame=lines,
	literate=
	*{0}{{{\color{numb}0}}}{1}
	{1}{{{\color{numb}1}}}{1}
	{2}{{{\color{numb}2}}}{1}
	{3}{{{\color{numb}3}}}{1}
	{4}{{{\color{numb}4}}}{1}
	{5}{{{\color{numb}5}}}{1}
	{6}{{{\color{numb}6}}}{1}
	{7}{{{\color{numb}7}}}{1}
	{8}{{{\color{numb}8}}}{1}
	{9}{{{\color{numb}9}}}{1}
	{:}{{{\color{punct}{:}}}}{1}
	{,}{{{\color{punct}{,}}}}{1}
	{\{}{{{\color{delim}{\{}}}}{1}
	{\}}{{{\color{delim}{\}}}}}{1}
	{[}{{{\color{delim}{[}}}}{1}
	{]}{{{\color{delim}{]}}}}{1},
}

\usepackage{graphicx}
\usepackage{tcolorbox}


\newtcolorbox{remarkbox}[1][]{
	colback=yellow!10,
	colframe=red!70!black,
	title=Definicja,
	fonttitle=\bfseries,
	boxrule=1pt,
	arc=4pt,
	leftrule=2pt,
	#1
}
\newtcolorbox{infobox}[1][]{
	colback=yellow!10,
	colframe=red!70!black,
	fonttitle=\huge\bfseries,
	fontupper=\Huge\color{red},
	boxrule=1pt,
	arc=4pt,
	leftrule=2pt,
	top=10mm,
	bottom=10mm,
	left=10mm,
	right=10mm,
	center title,        % wycentrowanie tytułu
	halign=center,       % wycentrowanie tekstu poziomo
	valign=center,       % wycentrowanie tekstu pionowo
	height=\textheight,  % wysokość boxa = wysokość strony
	width=\textwidth,    % szerokość boxa = szerokość strony
	before skip=0pt,     % brak dodatkowego odstępu przed
	after skip=0pt,      % brak dodatkowego odstępu po
	#1
}
\graphicspath{ {./img/} }
\newcounter{zadanieCounter}

\stepcounter{zadanieCounter}
\newcommand{\Z}{Zadanie \thezadanieCounter \stepcounter{zadanieCounter}}
\newcommand{\sectionZadanie}{\section*{\Z}}
\newcommand{\sectionZadania}{\section*{Zadania}}
\definecolor{codegreen}{rgb}{0,0.6,0}
\definecolor{codegray}{rgb}{0.5,0.5,0.5}
\definecolor{codepurple}{rgb}{0.58,0,0.82}
\definecolor{backcolour}{rgb}{0.95,0.95,0.92}
\lstdefinestyle{mystyle}{
	backgroundcolor=\color{backcolour},   
	commentstyle=\color{codegreen},
	keywordstyle=\color{magenta},
	numberstyle=\tiny\color{codegray},
	stringstyle=\color{codepurple},
	basicstyle=\ttfamily\footnotesize,
	breakatwhitespace=false,         
	breaklines=true,                 
	captionpos=b,                    
	keepspaces=true,                 
	numbers=left,                    
	numbersep=5pt,                  
	showspaces=false,                
	showstringspaces=false,
	showtabs=false,                  
	tabsize=4
}

%opening
\title{\pdfTitle}
\author{
	1 \\ e-mail  \href{1}{1} \and
	2 \\ e-mail  \href{2}{2}
}
\lstset{style=mystyle}

\usepackage{xcolor}
\usepackage{tcolorbox}
\tcbuselibrary{minted} % umożliwia użycie minted w tcolorbox
\newtcblisting{CodeBox}[2][]{
	listing engine=minted,
	minted language=#2,
	minted options={
		linenos,
		breaklines,
		fontsize=\Large
	},
	colback=black!3,
	colframe=black!70,
	title=#1,
	sharp corners,
	boxrule=0.8pt,
	left=6pt,
	right=6pt,
	top=6pt,
	bottom=6pt
}
\newtcblisting{CodeBoxJsx}[1][]{
	listing engine=minted,
	minted language=jsx,
	minted options={
		linenos,
		breaklines,
		fontsize=\Large
	},
	colback=black!3,
	colframe=black!70,
	title=#1,
	sharp corners,
	boxrule=0.8pt,
	left=6pt,
	right=6pt,
	top=6pt,
	bottom=6pt
}
\definecolor{colorLazoYellow1Light}{RGB}{255,248,230}
\definecolor{colorLazoBlue1Light}{RGB}{227,243,255}
\newtcblisting{jscode}{colback=colorLazoBlue1Light,listing only,listing engine=minted,minted language=javascript, 	minted options={
		linenos,
		breaklines
	},
	sharp corners,
	boxrule=0.8pt,
	left=6pt,
	right=6pt,
	top=6pt,
	bottom=6pt
}
\newtcblisting{jsxcode}{colback=colorLazoBlue1Light,listing only,listing engine=minted,minted language=jsx, 	minted options={
		linenos,
		breaklines
	},
	sharp corners,
	boxrule=0.8pt,
	left=6pt,
	right=6pt,
	top=6pt,
	bottom=6pt
}

\newcommand{\MyCopyrightYear}{2024-\the\year}

\usepackage{luacode}


\begin{luacode*}
	local lfs = require("lfs")  -- LuaFileSystem
	
	function include_all_tex(folder)
	local files = {}
	for file in lfs.dir(folder) do
	if file:match("%.tex$") then
	table.insert(files, file)
	end
	end
	table.sort(files)  -- sortowanie alfabetyczne
	
	for _, file in ipairs(files) do
	-- LaTeX wymaga slasha jako separatora
	local path = folder:gsub("\\","/") .. file
	tex.print("\\input{" .. path .. "}")
	end
	end

\end{luacode*}

\begin{luacode}
	function month_year_nominative_pl()
	local months = {
		"Styczeń","Luty","Marzec","Kwiecień","Maj","Czerwiec",
		"Lipiec","Sierpień","Wrzesień","Październik","Listopad","Grudzień"
	}
	local t = os.date("*t")
	tex.print(months[t.month] .. " " .. t.year)
	end
\end{luacode}

\newcommand{\IncludeAllTex}[1]{%
	\directlua{include_all_tex("#1")}%
}
\newcommand{\monthyearpl}{\protect\directlua{month_year_nominative_pl()}}